% ******************************* Plantilla para informes de laboratorio para la Escuela de Ciencia e Ingenieria en Materiales **************************
%Tecnológico de Costa Rica
\documentclass[letterpaper]{article}


%%%%%%%%%%%%
%Paquetes
%%%%%%%%%%%%
%\documentclass[journal=esthag,manuscript=article]{achemso}
\usepackage[spanish]{babel}
\usepackage[utf8]{inputenc}
\usepackage{fancyhdr}
\usepackage{authblk}
\usepackage{amsfonts}
\usepackage[version=3]{mhchem}
\usepackage{textgreek}
\usepackage{pdflscape} % landscape pages
\usepackage{lmodern} % Previene advertencias sobre fuentes no disponibles
\usepackage{graphicx}
\usepackage[left=3cm,right=2cm,top=2cm,bottom=3cm]{geometry} % Configura los margenes de página
%\usepackage{lineno} %Numeros en las lineas
%\linenumbers
%Paquete para incluir hipervínculos
\usepackage[colorlinks=true, 
linkcolor = black,
urlcolor  = blue,
citecolor = black,
anchorcolor = blue]{hyperref}
\usepackage{lipsum}
\usepackage{Sweave}
\usepackage{multirow}
\usepackage[table,xcdraw]{xcolor}
\usepackage{array}
\usepackage{tabu}
\usepackage{caption} %Permite escribir en negrita los titulos de las figuras y las tablas.
\usepackage[labelfont=bf]{caption} %Configura los titulos de las figuras y las tablas en negrita.

%Paquete para el manejo de las unidades en el Sistema Internacional
\usepackage{siunitx}
\sisetup{mode=text, output-decimal-marker = {,}, per-mode = symbol, qualifier-mode = phrase, qualifier-phrase = { de }, list-units = brackets, range-units = brackets, range-phrase = --}
\DeclareSIUnit[number-unit-product = \;] \atmosphere{atm}
\DeclareSIUnit[number-unit-product = \;] \pound{lb}
\DeclareSIUnit[number-unit-product = \;] \inch{"}
\DeclareSIUnit[number-unit-product = \;] \foot{ft}
\DeclareSIUnit[number-unit-product = \;] \yard{yd}
\DeclareSIUnit[number-unit-product = \;] \mile{mi}
\DeclareSIUnit[number-unit-product = \;] \pint{pt}
\DeclareSIUnit[number-unit-product = \;] \quart{qt}
\DeclareSIUnit[number-unit-product = \;] \flounce{fl-oz}
\DeclareSIUnit[number-unit-product = \;] \ounce{oz}
\DeclareSIUnit[number-unit-product = \;] \degreeFahrenheit{\SIUnitSymbolDegree F}
\DeclareSIUnit[number-unit-product = \;] \degreeRankine{\SIUnitSymbolDegree R}
\DeclareSIUnit[number-unit-product = \;] \usgallon{galón}
\DeclareSIUnit[number-unit-product = \;] \uma{uma}
\DeclareSIUnit[number-unit-product = \;] \ppm{ppm}
\DeclareSIUnit[number-unit-product = \;] \eqg{eq-g}
\DeclareSIUnit[number-unit-product = \;] \normal{\eqg\per\liter\of{solución}}
\DeclareSIUnit[number-unit-product = \;] \molal{\mole\per\kilo\gram\of{solvente}}


\addto\captionsspanish{\renewcommand{\tablename}{Tabla}} %Cambia el nombre de cuadros a tablas
%%%%%%%%%%%%%
%Encabezado
%%%%%%%%%%%%%  
\title{\textbf{Título del Laboratorio Realizado}}
\author[1]{Primer Integrante}
\author[2]{Segundo Integrante}
\author[3]{Tercer Integrante}
\author[4]{Cuarto Integrante}
\author[5]{Quinto Integrante}
\affil[1-5]{Escuela de Ciencia e Ingeniería en Materiales, Tecnológico de Costa Rica}
\affil[*]{Este trabajo evalua el atributo de investigacion nivel medio.}


%Imagenes en el encabezado de página
\lhead{\begin{picture}(0,0) \put(0,0){\includegraphics[width=30mm]{./images/ECIM.jpg}} \end{picture}}
\rhead{\begin{picture}(0,0) \put(-86.7,0){\includegraphics[width=30mm]{./images/ima3.jpg}} \end{picture}}
\renewcommand{\headrulewidth}{0.5pt}
\pagestyle{fancy}

%%%%%%%%%%%%%%%%%%%%%%%
%Inicio del documento
%%%%%%%%%%%%%%%%%%%%%%%
\begin{document}
	
	
%\input{2015-05-30_PEEC_trip_report-concordance}
\maketitle
\renewenvironment{abstract}
 {\quotation\small\noindent\rule{\linewidth}{.5pt}\par\smallskip
  {\centering\bfseries\abstractname\par}\medskip}
 {\par\noindent\rule{\linewidth}{.5pt}\endquotation}

%%%%%%%%%%
%Resumen
%%%%%%%%%%
\begin{abstract}
El resumen es una versión condensada de todo el informe escrita en apenas 100
a 200 palabras. Para lograr este propósito el resumen debe contener una o dos frases
introductorias al tema, la descripción del problema a resolver, el como fue resuelto,
la metodología empleada, los principales resultados y/o conclusiones del trabajo. En
la practica es mas fácil escribir el resumen al final, cuando las demás secciones estén
listas.
\end{abstract}

%%%%%%%%%%%%%%%
%Introducción
%%%%%%%%%%%%%%%
\section{Introducción}

Esta sección debe incluir información general que explique la importancia del experimento. Se deben utilizar referencias para soportar la base científica del mismo; esto le da credibilidad, porque indica que preparó adecuadamente la práctica. Debe incluir la información relevante relacionada con los experimentos; haga especial énfasis en condiciones inusuales o críticas. Las referencias deben ir numeradas consecutivamente en el texto y aparecer como una lista en la sección de referencias \cite{Cooper2009}. Las referencias más comunes incluyen libros, manuales de laboratorio y artículos.

%%%%%%%%%%%%%%%%%%%%%%%
%Materiales y Métodos
%%%%%%%%%%%%%%%%%%%%%%%
\section{Materiales y Métodos}

En esta sección se debe presentar de forma condensada los datos obtenidos durante el experimento. Todos los datos deben ser resumidos en tablas y/o gráficas enumeradas. El título debe contener una descripción clara del contenido. 
Las tablas deben reservarse para reportar listas de valores, mientras que la relaciones entre variables deben reportarse en gráficas. Tanto las gráficas, los esquemas de reacción, las fotografías y las ilustraciones deben agruparse bajo el título figura. En la figura \ref{fig:esquema} se observa un ejemplo de un esquema de reacción.

\begin{figure}[h!]
	\centering
	\includegraphics[width=0.7\linewidth]{images/esquema.eps}
	\caption{Esquema de reacción}
	\label{fig:esquema}
\end{figure}

\begin{table}[h!]
	\centering
	\begin{tabu} to 0.9\textwidth {  X[c]  X[c]  X[c]  }
		\hline
		\textbf{Papel Indicador}               & \textbf{Viraje de Color}            & \textbf{Rango de pH}     \\ \hline
		\multicolumn{1}{l}{\multirow{2}{*}{Papel tornasol neutro}} & violeta $\rightarrow$ rojo          & \num{7.0} -- \num{5.0}   \\
											   & violeta $\rightarrow$ azul          & \num{7.0} -- \num{8.0}   \\
		\multicolumn{1}{l}{Papel tornasol azul}                    & azul $\rightarrow$ rojo             & \num{8.0} -- \num{5.0}   \\
		\multicolumn{1}{l}{Papel tornasol rojo}                    & rojo $\rightarrow$ azul             & \num{5.0} -- \num{8.0}   \\
		\multicolumn{1}{l}{Papel amarillo brillante}			   & amarillo $\rightarrow$ rojo         & \num{6.7} -- \num{7.9}   \\
		\multicolumn{1}{l}{Papel Congo}							   & azul $\rightarrow$ rojo             & \num{8.0} -- \num{5.0}   \\
		\multicolumn{1}{l}{Papel amarillo nitrazina}			   & amarillo $\rightarrow$ violeta-azul & \num{6.0} -- \num{7.0}   \\
		\multicolumn{1}{l}{Papel fenolftaleína}					   & blanco $\rightarrow$ rojo           & \num{8.3} -- \num{10.0}  \\ \hline
	\end{tabu}
	\caption{Papeles indicadores y cambios de pH \cite{Lozano2010}.}
	\label{tab:indicadorespH}
\end{table}

\begin{figure}[h!]
		\centering
		\includegraphics[width=10cm]{./images/tubes.jpg}
		\caption{Resultados del ensayo de insaturación}
		\label{fig:tubes}
\end{figure}

\begin{table}[h!]
	\centering
	\begin{tabu} to 0.9\textwidth {  X[c]  X[c]  X[c] X[c] }
		\hline
		\textbf{Compuesto} & \textbf{Punto de ebullición experimental \SI{}{\degreeCelsius}} & \textbf{Punto de ebullición corregido \SI{}{\degreeCelsius}} & \textbf{Punto de ebullición literatura \SI{}{\degreeCelsius}} \\ \hline
		Tolueno            & b                                         & c                                      & \num{110.6}                                   \\
		Ciclohexano        & b                                         & c                                      & \num{80,74}                                   \\
		Ciclohexanol       & b                                         & \num{110.6}                                  & \num{160,8}                                   \\ \hline
	\end{tabu}
	\caption{Puntos de ebullición de las muestras problema.}
	\label{tab:boiling}
\end{table}

\subsection{Muestra de cálculos}

\begin{equation}
E={ E }^{ 0 }-\frac { 0.05916 }{ 2 } \log _{ 10 }{ \frac { \frac { { \left[ { H }^{ + } \right]  }^{ 2 }{ F }_{ { H }_{ 2 }A } }{ { \left[ { H }^{ + } \right]  }^{ 2 }+{ \left[ { H }^{ + } \right]  }{ K }_{ 1 }+{ K }_{ 1 }{ K }_{ 2 } }  }{ { F }_{ D }{ \left[ { H }^{ + } \right]  }^{ 2 } }  } 
\end{equation}

%%%%%%%%%%%%%%%%%%%%%%%%
%Resultados y Análisis
%%%%%%%%%%%%%%%%%%%%%%%%
\section{Resultados y Análisis}

\subsection{Corrección de los puntos de ebullición}

Los puntos de ebullición fueron corregidos utilizando la ecuación de Sydney-Young \cite{young:1902} (Ecuación \ref{eq:bpcorrection}). En donde \textit{t} es el punto de ebullición experimental, $C = \dfrac{dP}{dt} \times \dfrac{1}{T}$ es un valor aproximadamente constante que puede aproximarse a \num{0.00012} para líquidos polares y a \num{0.00010} para líquidos no polares y \textit{p} es la presión atmosférica del laboratorio expresada en mmHg.

\begin{equation}
\label{eq:bpcorrection}
\Delta T = C(760 - p)(273 + t)
\end{equation}

Se encontraron desviaciones respecto a los valores experimentales reportados en la literatura como se observa en la tabla \ref{tab:boiling}. La diferencia observada se le puede atribuir al menos a tres razones; primero la presión atmosférica del laboratorio en el momento del experimento no se pudo medir así que se tomó un valor de \SI{680}{\mmHg}. Segundo el termómetro que se utilizó no fue calibrado previamente además que era un termómetro de inmersión completa y en el experimento solo estaba sumergido el bulbo, por lo que no existe certeza sobre la exactitud de la medida. Finalmente, el valor hay que recordar que los valores de la constante \textit{C} son aproximados y que en realidad esta constante es diferente para cada líquido.

\subsection{Ecuaciones Químicas}

A continuación se muestra un ejemplo de una reacción química escrita con la ayuda del paquete \textit{mhchem}.

\begin{center}
	\ce{1/2 O2 _{(g)} + 2 H+ + 2 e^- <=> H2O} \hspace{19mm} $\times 3$ \hspace{20mm} $E^0 = \SI{1.229}{\volt}$\\
	\vspace{3mm}
	\ce{2 Cr^{3+} + 7 H2O <=> Cr2O7^{2-} + 14 H+ + 6 e^-} \hspace{24mm} $E^0 = \SI{-1.36}{\volt}$\\
	\underline{\hspace{145mm}}
	
	\ce{3 O2 _{(g)} + 4 Cr^{3+} + 8 H2O <=> 2 Cr2O7^{2-} + 16 H+ } \hspace{15mm} $E^0 = \SI{-0.13}{\volt}$\\
\end{center}

%%%%%%%%%%%%%%%
%Conclusiones
%%%%%%%%%%%%%%%
\section{Conclusiones}
\num{3.1415592d-3}

\begin{equation}
\label{eq:ensayo}
x=3y+5
\end{equation}

Ahora quiero referenciar la ecuación \ref{eq:ensayo} y continuo escribiendo

\ce{H2SO4 + Ca(OH)2 -> CaSO4 v + 2 H2O ^}

\begin{equation}
\label{eq:cuadratica}
x=\frac { -b\pm \sqrt [ 2 ]{ { b }^{ 2 }-4ac }  }{ 2a } 
\end{equation}

Si necesito referirme a la ecuación cuadrática (ver ecuación \ref{eq:cuadratica}) y continúo escribiendo

%%%%%%%%%%%%%%%
% Referencias
%%%%%%%%%%%%%%%
\bibliographystyle{plain}
\bibliography{PEECbib}

\newpage
%%%%%%%%%
%Anexos
%%%%%%%%%
\section{Anexos}
%En caso de utilizar anexos se colocaran debajo de esta línea




%%%%%%%%%%%%%%%%%%%%%%%%%%%%%%%%%%%%%%%%%%%%%%%%%%%%%%%%%%%%%%%%%%%%%%%%%%%%%%%%%%%%%%%%%%%%%%%
%No modificar esta sección
\begin{table}[h!]
	\centering
	\begin{tabu} to 0.9\textwidth {  X[c]  X[c]  X[c]  }
		\hline
		\textbf{Item a evaluar} & \textbf{Valor} & \textbf{Valor Obtenido} \\ \hline
		Presentacion            & 5 \%             &                         \\ \hline
		Resumen                 & 10 \%            &                         \\ \hline
		Introduccion            & 10 \%            &                         \\ \hline
		Materiales y Metodos    & 15 \%            &                         \\ \hline
		Resultados y Discusion  & 40 \%            &                         \\ \hline
		Conclusiones            & 15 \%            &                         \\ \hline
		Bibliografia            & 5 \%             &                         \\ \hline
	\end{tabu}
	\caption{Rubrica de evaluación.}
	\label{Rubros}
\end{table}

\newpage
\begin{table}[h!]
	\centering
	\begin{tabu} to 1\textwidth { | X[c] | X[c] | X[c] | X[c] | X[c] | }
		\hline
		\multirow{3}{*}{Dimensión}                                            & \multicolumn{4}{c|}{Niveles de desempeńo}                                                                                                                                                                                                                                                                                                                                                                                                                                                                                                                                      \\ \cline{2-5} 
		& 1                                                                                                                                                       & 2                                                                                                               & 3                                                                                                                                        & 4                                                                                                                                                       \\ \cline{2-5} 
		& Insatisfactorio                                                                                                                                         & Necesita Mejorar                                                                                                & Satisfactorio                                                                                                                            & Sobresaliente                                                                                                                                           \\ \hline
		Identificación de problemas complejos de ingenería.                   & Identifica un problema a nivel elemental en ingeniería.                                                                                                 & Identifica problemas complejos de ingeniería en función de variables, pero sin considerar las causas y efectos. & Identifica problemas complejos* de ingeniería en función de variables, causas y efectos.                                                 & Identifica problemas complejos de ingeniería en función de variables, causas, efectos, impacto e importancia.\vspace{5pt}                                           \\ \hline
		Estado del arte.                                                      & Ausencia de síntesis de investigaciones pertinentes al problema que fundamenta el marco teórico con  referentes bibliográficos actualizados.             & Realiza una síntesis de investigaciones pertinentes al problema, sin una fundamentación del marco teórico.      & Realiza síntesis de investigaciones pertinentes al problema que fundamenta el marco teórico con  referentes bibliográficos actualizados. & Realiza síntesis, integra y compara investigaciones pertinentes al problema que fundamenta el marco teórico con  referentes bibliográficos actualizados.\vspace{5pt} \\ \hline
		Diseño de metodologías                                                & Carece de una definición del tipo de estudio, población, muestra,  instrumentos y técnicas para la recolección de datos y el procedimiento de análisis. & Define el tipo de estudio, población, muestra,  instrumentos y técnicas, sin coherencia entre los elementos.    & Define el tipo de estudio, población, muestra,  instrumentos y técnicas para la recolección de datos y el procedimiento de análisis.     & Define y argumenta el tipo de estudio, población, muestra,  instrumentos y técnicas para la recolección de datos y el procedimiento de análisis.\vspace{5pt}        \\ \hline
		Aplicación y prueba de experimentos (hipótesis o posibles soluciones) & Carece de  pruebas  (hipótesis o posibles soluciones), según los  parámetros  establecidos en el  diseño experimental.                                  & Realiza diversas pruebas, sin considerar los  parámetros establecidos en el  experimento.                       & Realiza  pruebas  (hipótesis o posibles soluciones), según los  parámetros  establecidos en el  diseño experimental.                     & Realiza  pruebas  (hipótesis o posibles soluciones), según los  parámetros  establecidos en el  diseño experimental y compara los resultados.\vspace{5pt}           \\ \hline
	\end{tabu}
	\caption{Tabla de atributos Investigación}
	\label{IN}
\end{table}
%%%%%%%%%%%%%%%%%%%%%%%%%%%%%%%%%%%%%%%%%%%%%%%%%%%%%%%%%%%%%%%%%%%%%%%%%%%%%%%%%%%%%%%%%%%%%%%
\end{document}
